% GENERATED FILE. Edit canonical lessons, then run npm run generate:books.
% canonical-source-hash: fnv1a64:6f80b5993b6e50ae
% canonical-lessons: KA-C256-tollu, KA-C256-cappate, KA-C256-donku, KA-C256-masuku, KA-C256-kikkirida

\chapter{Hollow, Flat, Crooked, Blurred, Crowded}
\label{ch:ka-hollow-flat-crooked-blurred-crowded}
\hlchaptermodality{256}

\begin{chapteropening}
\textbf{By the end of this chapter:} \emph{I can say hollow, flat, crooked, blurred and crowded.}

Complete the last lesson of chapter 256: I can say hollow, flat, crooked, blurred and crowded.
\end{chapteropening}

\section[\texorpdfstring{\kn{ಟೊಳ್ಳು} (ṭoḷḷu)}{ಟೊಳ್ಳು (ṭoḷḷu)}]{\kn{ಟೊಳ್ಳು} (ṭoḷḷu) — hollow}

\label{lesson:KA-C256-tollu}



\begin{quote}
Before the new one: say the Kannada for serious, then the Kannada for round.
\end{quote}

\subsection*{You'll want to know: \kn{ಟೊಳ್ಳು}}
\textbf{\kn{ಟೊಳ್ಳು}} — \emph{ṭoḷḷu} — \textquotedblleft{}hollow\textquotedblright{}.

\begin{grammarlens}[title={What you've built}]
One more word for this chapter.
\end{grammarlens}

\subsection*{Your turn}
\begin{itemize}
\raggedright
  \item \emph{Say it:} \emph{ṭoḷḷu}
  \item \emph{Say it:} \emph{ṭoḷḷu}, once more
  \item \emph{Say it:} say \emph{ṭoḷḷu}
  \item \emph{Recall:} say the Kannada for serious, then the Kannada for round, then say \emph{ṭoḷḷu} again
\end{itemize}

\begin{tcolorbox}[breakable,colback=teal!4,colframe=teal!35!black,title={Before you move on}]
What does \kn{ಟೊಳ್ಳು} mean? (\textquotedblleft{}Hollow\textquotedblright{}.) Say it once more.
\end{tcolorbox}
\section[\texorpdfstring{\kn{ಚಪ್ಪಟೆ} (cappaṭe)}{ಚಪ್ಪಟೆ (cappaṭe)}]{\kn{ಚಪ್ಪಟೆ} (cappaṭe) — flat}

\label{lesson:KA-C256-cappate}



\begin{quote}
Before the new one: say the Kannada for round, then the Kannada for hollow.
\end{quote}

\subsection*{You'll want to know: \kn{ಚಪ್ಪಟೆ}}
\textbf{\kn{ಚಪ್ಪಟೆ}} — \emph{cappaṭe} — \textquotedblleft{}flat\textquotedblright{}.

\begin{grammarlens}[title={What you've built}]
One more word for this chapter.
\end{grammarlens}

\subsection*{Your turn}
\begin{itemize}
\raggedright
  \item \emph{Say it:} \emph{cappaṭe}
  \item \emph{Say it:} \emph{cappaṭe}, once more
  \item \emph{Say it:} say \emph{cappaṭe}
  \item \emph{Recall:} say the Kannada for round, then the Kannada for hollow, then say \emph{cappaṭe} again
\end{itemize}

\begin{tcolorbox}[breakable,colback=teal!4,colframe=teal!35!black,title={Before you move on}]
What does \kn{ಚಪ್ಪಟೆ} mean? (\textquotedblleft{}Flat\textquotedblright{}.) Say it once more.
\end{tcolorbox}
\section[\texorpdfstring{\kn{ಡೊಂಕು} (ḍoṅku)}{ಡೊಂಕು (ḍoṅku)}]{\kn{ಡೊಂಕು} (ḍoṅku) — crooked, bent}

\label{lesson:KA-C256-donku}



\begin{quote}
Before the new one: say the Kannada for hollow, then the Kannada for flat.
\end{quote}

\subsection*{You'll want to know: \kn{ಡೊಂಕು}}
\textbf{\kn{ಡೊಂಕು}} — \emph{ḍoṅku} — \textquotedblleft{}crooked, bent\textquotedblright{}.

\begin{grammarlens}[title={What you've built}]
One more word for this chapter.
\end{grammarlens}

\subsection*{Your turn}
\begin{itemize}
\raggedright
  \item \emph{Say it:} \emph{ḍoṅku}
  \item \emph{Say it:} \emph{ḍoṅku}, once more
  \item \emph{Say it:} say \emph{ḍoṅku}
  \item \emph{Recall:} say the Kannada for hollow, then the Kannada for flat, then say \emph{ḍoṅku} again
\end{itemize}

\begin{tcolorbox}[breakable,colback=teal!4,colframe=teal!35!black,title={Before you move on}]
What does \kn{ಡೊಂಕು} mean? (\textquotedblleft{}Crooked, bent\textquotedblright{}.) Say it once more.
\end{tcolorbox}
\section[\texorpdfstring{\kn{ಮಸುಕು} (masuku)}{ಮಸುಕು (masuku)}]{\kn{ಮಸುಕು} (masuku) — blurred, dim}

\label{lesson:KA-C256-masuku}



\begin{quote}
Before the new one: say the Kannada for flat, then the Kannada for crooked.
\end{quote}

\subsection*{You'll want to know: \kn{ಮಸುಕು}}
\textbf{\kn{ಮಸುಕು}} — \emph{masuku} — \textquotedblleft{}blurred, dim\textquotedblright{}.

\begin{grammarlens}[title={What you've built}]
One more word for this chapter.
\end{grammarlens}

\subsection*{Your turn}
\begin{itemize}
\raggedright
  \item \emph{Say it:} \emph{masuku}
  \item \emph{Say it:} \emph{masuku}, once more
  \item \emph{Say it:} say \emph{masuku}
  \item \emph{Recall:} say the Kannada for flat, then the Kannada for crooked, then say \emph{masuku} again
\end{itemize}

\begin{tcolorbox}[breakable,colback=teal!4,colframe=teal!35!black,title={Before you move on}]
What does \kn{ಮಸುಕು} mean? (\textquotedblleft{}Blurred, dim\textquotedblright{}.) Say it once more.
\end{tcolorbox}
\section[\texorpdfstring{\kn{ಕಿಕ್ಕಿರಿದ} (kikkirida)}{ಕಿಕ್ಕಿರಿದ (kikkirida)}]{\kn{ಕಿಕ್ಕಿರಿದ} (kikkirida) — crowded, packed}

\label{lesson:KA-C256-kikkirida}



\begin{quote}
Before the new one: say the Kannada for crooked, then the Kannada for blurred.
\end{quote}

\subsection*{You'll want to know: \kn{ಕಿಕ್ಕಿರಿದ}}
\textbf{\kn{ಕಿಕ್ಕಿರಿದ}} — \emph{kikkirida} — \textquotedblleft{}crowded, packed\textquotedblright{}.

\begin{grammarlens}[title={What you've built}]
One more word for this chapter.
\end{grammarlens}

\subsection*{Your turn}
\begin{itemize}
\raggedright
  \item \emph{Say it:} \emph{kikkirida}
  \item \emph{Say it:} \emph{kikkirida}, once more
  \item \emph{Say it:} say \emph{kikkirida}
  \item \emph{Recall:} say the Kannada for crooked, then the Kannada for blurred, then say \emph{kikkirida} again
\end{itemize}

\begin{tcolorbox}[breakable,colback=teal!4,colframe=teal!35!black,title={Before you move on}]
What does \kn{ಕಿಕ್ಕಿರಿದ} mean? (\textquotedblleft{}Crowded, packed\textquotedblright{}.) Say it once more.
\end{tcolorbox}
